\section{Positive selection: a valid target and a finite-model gain}
\label{selection:section}

This section distinguishes two positive-selection targets. Vanishing total
loss and a uniform bound for the full retained collision law would force the
original joint law into \(L^2\), which is false for some prescribed Frostman
probabilities. Retaining a fixed positive mass is sufficient for the desired
distance conclusion. Persistent deletion is one possible implementation;
nesting is not required. We then prove a sharper transverse estimate for
the finite train-track test family. No improved dimension-only distance
criterion follows without an additional geometric argument.

\subsection{The compactness distinction}
\label{selection:compactness-section}
Let \(\mu,\nu\) be compactly supported planar probabilities and
\(\varphi\ge0\) a smooth compactly supported probability on \(\R\).
Write \(\varphi_\epsilon(t)=\epsilon^{-1}\varphi(t/\epsilon)\) and
\[
 \Lambda(dy,dt)=\nu(dy)(d_y)_*\mu(dt).
\]
A positive interaction selection is a measurable kernel
\(\eta_{j,y}\le\mu\). For \(\epsilon_j\downarrow0\), put
\[
 F_j(y,t)=\varphi_{\epsilon_j}*(d_y)_*\eta_{j,y}(t).
\]
All radial supports lie in one fixed bounded interval for sufficiently
large \(j\). Fix an exponent \(1<q<\infty\), independently of \(j\).

\begin{proposition}[Positive-mass compactness]
\label{selection:compactness}
Suppose \(\sup_j\|F_j\|_{L^q(d\nu\,dt)}<\infty\).
If
\[
 \int\eta_{j,y}(\R^2)\dd\nu(y)\ge m_0>0,
\]
then \(\Lambda\) has an absolutely continuous submeasure of mass at
least \(m_0\), with density in \(L^q(d\nu\,dt)\).
Consequently a positive-\(\nu\) family of pins has positive-length
distances from \(\supp\mu\).

If, more strongly,
\[
 \int\bigl(1-\eta_{j,y}(\R^2)\bigr)\dd\nu(y)\longrightarrow0,
\]
then the \emph{entire} joint law \(\Lambda\) has an \(L^q(d\nu\,dt)\)
density.
\end{proposition}
\begin{proof}
Pass to a weakly convergent subsequence in the reflexive space
\(L^q(d\nu\,dt)\), with limit \(F\). Positivity passes to the limit.
All bounded test functions used below belong to the dual space on the fixed
finite joint support.
Testing against the indicator of an interval containing all radial supports
shows \(\int F\dd\nu\,dt\ge m_0\).
The selected mollified laws are dominated by the full mollifications of
\(\Lambda\). Testing against any nonnegative continuous compactly
supported function and taking the limit gives
\[
 F\dd\nu\,dt\le\Lambda.
\]
Disintegration gives \(F(y,t)dt\le(d_y)_*\mu\) for \(\nu\)-almost
every \(y\). It is nonzero on a positive-\(\nu\) set of pins and
supported on the corresponding compact distance sets. Those sets therefore
have positive Lebesgue measure.

Under the stronger hypothesis, the total variation difference between the
selected and full mollified joint laws is at most the averaged deleted mass:
convolution and pushforward contract total variation. This tends to zero.
Full mollifications converge weakly to \(\Lambda\), so the same test
functions identify \(F\dd\nu\,dt=\Lambda\).
\end{proof}

\begin{remark}[A false dimension-only extension]
\label{selection:false-candidate}
For \(1<s<4/3\) and \(3-2/s\le u\le2\), the detailed companion
\href{https://github.com/CoolRmal/falconer-packing/blob/main/docs/packing-unforced/2026-09-30-off-cluster-collision-obstruction.md}
{off-cluster collision note} constructs fixed separated probabilities with
uniform \(s\)-Frostman and \(u\)-covering bounds whose fully mollified
joint squared \(L^2\) norms diverge along a sequence. Their support has
Hausdorff and packing dimensions both \(s\). A variant also has finite
energy at the stated Frostman exponent.

The construction uses scaled train-track blocks with
\[
 \lambda_j=c2^{-j},\quad \delta_j=2^{-(j+j_0)^2},\quad
 w_j\asymp\lambda_j^s,\quad \epsilon_j\asymp\lambda_j\delta_j,
\]
and gives the positive collision lower bound
\[
 c\lambda_j^{3s-1}\delta_j^{3s/2-2}\longrightarrow\infty.
\]
The uniform Frostman, covering, gluing, and collision proofs are supplied
in that companion note; this finite construction is not reproduced here.

For this fixed pair, no positive selection with total deleted mass tending
to zero can have uniformly bounded full-law collision norms, even along an
arbitrary sequence of resolutions tending to zero. Otherwise
Proposition~\ref{selection:compactness}, with \(q=2\), would put the original joint law
in \(L^2\), and Young's convolution inequality would bound all its
mollified norms, contradicting the displayed divergence. Thus a candidate
asserting this conclusion for every prescribed pair is false for \emph{any}
vanishing deletion rate, not only a fast power rate. This is not an
obstruction to positive-length distances.
\end{remark}

The same distinction applies below exponent two. Section 5 of the
\href{https://github.com/CoolRmal/falconer-packing/blob/main/docs/packing-unforced/2026-09-30-subquadratic-profile-transfer.md}
{subquadratic transfer companion note} proves that the fixed train-track
joint law has no weak \(L^q\) density whenever
\[
 q>\frac{2-s}{3-2s},\qquad 1<s<4/3.
\]
It therefore cannot admit vanishing-loss selections uniformly bounded in
strong \(L^q\) for such a fixed exponent, by
Proposition~\ref{selection:compactness}. The cited construction and its
quantitative weak-norm proof are not reproduced here. This supplies no
obstruction for exponents sufficiently close to one.

A valid persistent implementation would choose discarded interaction sets
\(A_k\subset\supp\mu\times\supp\nu\) with
\[
 \sum_{k\ge k_0}(\mu\times\nu)(A_k)<1
\]
and retain
\[
 \eta_{j,y}=\mu\big|_{\{x:(x,y)\notin\bigcup_{k_0\le k\le j}A_k\}}.
\]
The mass stays uniformly positive. A uniform retained collision estimate
would then suffice by Proposition~\ref{selection:compactness}. Such an
estimate must use the accumulated selection: a uniform full-law bound after
only a current-scale deletion whose mass vanishes would have the stronger,
false consequence above.

For fixed measures, existence of a positive interaction restriction of
positive mass with an \(L^2\) joint distance density is equivalent to
\(\Lambda\) having a nonzero absolutely continuous component. For the
nontrivial direction, remove a null set carrying the singular part of
\(\Lambda\), truncate its remaining density to \(0<f\le M\), and
lift that restriction through \((y,x)\mapsto(y,|x-y|)\). A sufficiently
large finite \(M\) retains positive mass with a bounded density.
Allowing arbitrarily small \emph{fixed} loss similarly characterizes
absolute continuity; the \(L^2\) constant may then deteriorate with the
loss. The geometric construction of such a restriction is the substantive
missing step.

\subsection{The finite train-track model}
\label{selection:finite-model}
Fix
\[
 1<s<4/3,\quad 3-2/s\le u\le2,\quad
 \alpha=u-1,\quad\beta=(s-1)/2,\quad p=\beta/\alpha,
\]
and, for small \(\delta>0\), put
\[
 a=\delta^p,\qquad w=\delta^{1/2},\qquad h=\delta^{s/2},
 \qquad r_0=h/w=\delta^\beta.
\]
Thus \(0<\beta\le p<1/2\) and \(a\gg w\gg h\gg\delta\).
Take a maximal \(3a\)-separated net \(\{c_i\}_{i=1}^{N_T}\) in
an \(\alpha\)-Ahlfors regular subset of a fixed horizontal interval.
With \(m=N_T^{-1}\), standard disjoint-ball comparisons give
\begin{equation}\label{selection:center-count}
 m\asymp a^\alpha=\delta^\beta,\qquad
 \sum_{c_i\in J}m\le C(|J|^\alpha+a^\alpha).
\end{equation}
They also give a lower bound \(c\rho^\alpha\) for track mass within
radius \(\rho\) of any center when \(C_0a\le\rho\le c_0\), with
fixed constants. This follows by covering the underlying regular set with
net balls of radius \(3a\) and comparing its regular mass with the
upper mass of one ball.

In each track place \(N_S\asymp h^{-1}\) slats at heights \(kh\)
in a bounded interval. Each slat has \(N_X\asymp w/\delta\) disjoint
microscopic cells, with horizontal centers spaced by \(4\delta\) across
width comparable to \(w\). Each cell lies within \(C\delta\) of
its center and has mass \((N_TN_SN_X)^{-1}\asymp m h\delta/w\).
Inside a cell use any compact probability supported there. Let
\(\mu_\delta\) give equal mass to all cells and let \(\nu_\delta\)
be its translate by \((0,3)\), normalizing the vertical source interval
so source--pin distances are bounded above and below.
The regular microscopic tails giving the stated dimension and covering
properties of the original example are one permitted choice. The following
estimates only use the explicit geometry and counts just listed.

\begin{lemma}[One-track estimate]\label{selection:one-track}
Fix a pin in the translated source region and a track at horizontal offset
\(r=|c_i-y_1|\gg w\). Its mollified pinned density \(f_i\), at
resolution \(\epsilon\asymp\delta\), satisfies
\[
 \|f_i\|_1=m,\qquad
 \|f_i\|_\infty\le Cm(1+r_0/r),\qquad
 \|f_i\|_2\le Cm\sqrt{1+r_0/r}.
\]
Constants may depend on the fixed mollifier and resolution ratio.
\end{lemma}
\begin{proof}
Throughout the track the horizontal derivative of distance has constant
sign and magnitude comparable to \(r\). In one slat, a radial interval
of width \(C\delta\) therefore meets at most \(C/r\) cell centers.
The slat contributes density at most \(Cm h/(rw)\).
Its whole radial image, after mollification, has length \(O(rw+\delta)\),
which is \(O(rw)\) because \(r\gg w\) and \(w^2=\delta\).
Central radial values of successive slats are separated by at least
\(ch\), since vertical source--pin separation is bounded below.
Thus at most \(C(1+rw/h)\) slats contribute at one radial value.
Multiplication gives the asserted supremum bound, including when
\(rw<h\). Interpolation with the first norm proves the last bound.
\end{proof}

Minkowski now gives, for any retained family of tracks with \(r_i\gg w\),
\begin{equation}\label{selection:half-potential}
 \left\|\sum_i f_i\right\|_2
 \le C\left(1+\sqrt{r_0}\sum_i m r_i^{-1/2}\right).
\end{equation}
The source mass is at most one. Decomposing offsets into dyadic annuli and
using \eqref{selection:center-count}, for \(\rho\ge ca\), gives
\begin{equation}\label{selection:potential}
 \sum_{r_i\ge\rho}m r_i^{-1/2}\le
 \begin{cases}
 C_\alpha\rho^{\alpha-1/2},&\alpha<1/2,\\
 C\log(2/\rho),&\alpha=1/2,\\
 C_\alpha,&\alpha>1/2.
 \end{cases}
\end{equation}
Indeed each annulus contributes at most \(C(2^j\rho)^{\alpha-1/2}\),
and summing the geometric series gives these three cases.

\begin{proposition}[A sharper positive finite-model selection]
\label{selection:finite-theorem}
If \(u\ge4/3\), deleting only the pin's own source track gives
uniformly bounded mollified \(L^2\) densities, with lost mass
\(m\asymp\delta^\beta\). If \(u<4/3\), deleting tracks at offsets
less than a sufficiently large constant times
\[
 \rho_*=r_0^{1/(1-2\alpha)}
\]
gives uniformly bounded mollified \(L^2\) densities, with lost mass at most
\(C\delta^{\alpha\beta/(1-2\alpha)}\).
For the particular rule deleting only the own track, the transition
\(u=4/3\) is exact: below it the retained averaged squared \(L^2\)
norm diverges, including its off-cell contribution.
\end{proposition}
\begin{proof}
After own-track deletion, the other offsets are at least \(ca\gg w\).
For \(1/3\le\alpha<1/2\), the potential factor in
\eqref{selection:half-potential} is bounded because
\[
 \sqrt{r_0}a^{\alpha-1/2}
 =\delta^{\beta(3/2-1/(2\alpha))}\le1.
\]
There is no logarithmic loss at \(\alpha=1/3\). At \(\alpha=1/2\),
use boundedness of \(\sqrt{r_0}\log(2/a)\); for \(\alpha>1/2\),
use the constant bound in \eqref{selection:potential}.

When \(\alpha<1/3\), one has \(\rho_*\gg a\gg w\), and
\(\sqrt{r_0}\rho_*^{\alpha-1/2}=1\). The same estimates give the
norm bound, while \eqref{selection:center-count} gives the removed mass
\(C\rho_*^\alpha\).

For the converse concerning own-track deletion, Ahlfors regularity and
maximality of the net provide a different center at distance between
\(3a\) and \(Ca\) from each center. Indeed, for large fixed \(C\),
regular mass in the ball of radius \(Ca\) exceeds that in the ball
of radius \(6a\); a net point within \(3a\) of a point in the
difference is a suitable neighbor. A pin in the original track is therefore
at offset comparable to \(a\) from a retained neighboring track.
The latter's mollified radial support has length at most \(CN_Saw\).
Its mass is \(m\), so Cauchy--Schwarz and positivity give
\[
 \|f_y^{\rm retained}\|_2^2\ge c\frac{m^2h}{aw}
 \asymp\delta^{3\beta-p}.
\]
The exponent \(3\beta-p=\beta(3-1/\alpha)\) is negative exactly
when \(\alpha<1/3\). Finally the same-cell contribution is at most
\[
 C\delta^{-1}\sum_{A\text{ cell}}\mu_\delta(A)^2
 \lesssim\delta^{s-1}\longrightarrow0,
\]
since each cell has mass comparable to \(\delta^s\). The divergent
remainder therefore comes from distinct cells.
\end{proof}

\subsection{A necessary deletion budget below the transition}
\label{selection:budget-section}
There remains a gap between the sufficient budget above and a necessary
one. For \(\alpha<1/3\), any positive selection with averaged squared
\(L^2\) norm bounded by a fixed \(C\) must delete average mass at least
\begin{equation}\label{selection:budget}
 c_C\delta^{2\alpha\beta/(1-\alpha)}.
\end{equation}
Here is the proof. At offsets \(ca\le r_i\le\rho\), with \(\rho\gg a\),
the initial track mass is at least \(c\rho^\alpha\), and the number
of tracks is at most \(C\rho^\alpha/m\). If \(q_i\) is the retained
mass of track \(i\), its radial support has length at most
\(CN_Sr_iw\); hence its squared norm is at least
\(cr_0q_i^2/r_i\). Positivity, weighted Cauchy--Schwarz, and a total
deleted mass \(r_y\) at the pin give
\[
 \|f_y^{\rm retained}\|_2^2
 \ge c r_0m\rho^{-1-\alpha}(c\rho^\alpha-r_y)_+^2.
\]
Average and use Jensen. If the averaged deletion is at most
\((c/2)\rho^\alpha\), the averaged norm is at least
\(c'r_0m\rho^{\alpha-1}\). Choose
\(\rho=A_C(r_0m)^{1/(1-\alpha)}\), with \(A_C>0\) fixed and small
enough to exceed the proposed norm bound. This radius still satisfies
\(\rho\gg a\) for small \(\delta\), precisely because \(\alpha<1/3\).
The deletion must therefore exceed \((c/2)\rho^\alpha\), proving
\eqref{selection:budget}.

The sufficient and necessary exponents agree at \(\alpha=1/3\), but
not below it. Exact optimal deletion rates are not asserted there.
More substantially, the finite-model proof uses slat ordering, transverse
radial derivatives, and explicit counts absent from dimension data alone.
A dimension-only construction of positive-mass good interactions with a
corresponding retained collision estimate remains unproved. The compactness
criterion identifies a valid target; it does not supply that geometric input.
